\chapter{INTRODUCTION} % Main chapter title
\label{ChapterIntroduction} % For referencing the chapter elsewhere,
% use \ref{ChapterIntroduction}

Digital services have been growing exponentially and changed the way
people are approaching technologies, allowing for the easy access to
any apps, platforms, and the whole digital ecosystem. 
The expansion mentioned has also led to a lot of dependence on digital contractual agreements, 
such as Terms of Service (ToS), Privacy Policies, and 
End-User License Agreements (EULAs) which users have to accept prior to accessing services.

Although these agreements are important since they provide 
legal relationships between users and service providers, they are 
often very long and complicated. Because of this, the user finds it 
difficult to read and fully understand them before accepting the 
service, which raises concerns about whether or not the users 
are truly given informed consent and if transparency exists.
But on the bright side, the recent development of AI and 
Natural Language Processing\index{NLP} have made it possible to automatically 
analyse complex text without requiring prior knowledge of the content. 
Techniques like semantic embeddings\index{Semantic Embedding}, knowledge graphs\index{Knowledge Graph}, and
Retrieval-Augmented Generation\index{RAG} allow intelligent systems 
to extract useful insights from unstructured documents, 
while also preserving the meaning.

By applying the advancements mentioned above, within an AI-based framework 
to analyze legal agreements, the RiskWise initiative aims to generate 
intelligent and risk-aware insights, that help users better understand 
legal and other publicly available documents.
This chapter discusses the background, motivation, and problem formulation of the proposed system.

\section{Context and Background Study}
As the extent of digital services and online 
platforms is expanding at a rapid rate, users are 
becoming more and more bound by Terms of Service (ToS), 
Privacy Policies and End User License Agreements (EULAs) 
in order to gain access to apps/services. 
What these documents do, is that they serve as legally binding 
agreements between service providers and users 
outlining their right, responsibilities, limitations of 
liability, data usage policies, termination conditions and 
other contractual terms related to the use of service.  
Such agreements are necessary in order to make legal relationships 
between parties in digital ecosystems. But since the general user 
experience is often poor, their practical effectiveness is often reduced. 
Although these agreements are important, studies and user behaviour data 
show that most users either rarely read these documents 
completely or do not fully understand them.
The documents being too lengthy, 
having complicated legal language, etc. are the 
main reasons that contribute to this.

This general disengagement results in a great divide between user
informed consent and user understanding, that has implications for
the legitimacy of informed consent in the digital realm. 
Therefore, users could often agree to the terms and not even be 
aware of the possibilities for the complete collection of their 
data, unilateral changes to policies, subscription renewals, etc. 
This could lead to tremendous risks for their privacy, financial obligations, 
and access to services. 
New developments in Artificial Intelligence, specifically Large Language Models (LLMs)\index{LLM} and
Natural Language Processing (NLP)\index{NLP}, have proven tremendous scope for the 
automatic understanding of complex texts. On top of that, new techniques 
such as Retrieval-Augmented Generation (RAG)\index{RAG}, semantic embeddings\index{Semantic Embedding}, and
knowledge graph\index{Knowledge Graph} representations enable structured understanding from 
unstructured texts (for example, semantic meaning and relationship). 
These new developments are useful when it comes to creation of 
intelligent systems that can assist in the comprehension of legal documents.

The RiskWise framework, in this technological context, is presented here as an 
AI-driven solution to deal with legal documents, identify potentially 
risky clauses and present findings in a simplified and user-centric way. 
Using the hybrid form of 2 retrieval and reasoning mechanism, the idea is 
to create a link between the complexity of a contract and user 
understanding to guide decision making in digital interactions.

\section{Need Analysis}
A lot of practical, technical and societal reasons for the need of an
automated Terms and Conditions analysis system exist. First of all, 
modern digital ecosystems involve high levels of communication with a 
variety of online services, each one of which requires users to accept 
contractual agreements.
Time-consuming and requiring legal literacy that non-standard 
users lack, reading these items by hand results in
surface acceptance but not much consideration of the content.

Second, entities often include clauses granting wide discretion such
as the exclusive power to change clauses, the broad flexibility to
share data and the rights to cancel accounts. When no analytical
support is provided, users are exposed to unknowingly accepting
unsatisfactory terms within a contract that may threaten their
privacy, financial obligations to customers or their access to
services.

Thirdly, today’s automated applications are primarily focused upon 
providing a description of the work that the users have done and 
fail to interpret the risk they have detected. 
Summarization may be an instrument to reduce the difficulty of 
documents but does not necessarily check for clauses, introducing 
foreseeable potential risks (and why they might be associated). 
Consequently, users need systems which can not only minimize 
but also strategically identify, classify and
situate contractual risks\index{Risk Analysis}
(and why they might be associated).

Consequently, users need systems
which can not only minimize but also strategically identify, classify
and situate contractual risks. Furthermore, the widespread global
trend towards digital privacy and protection of consumers as well as
compliance requirements demands that we find mechanisms that allow us
to promote transparency to give the right of informed consent to
participants. It allows the system to identify material clauses,
categorize potential risks, and walks the user through an interactive
exploration of contractual content, and greatly increases their
knowledge and decision-making capacity. Thus, there is a great demand
for the proposed intelligent, user-friendly and privacy-preserving
system for the assistance of users to weigh legal contracts prior to 
acceptance.

\section{Problem Identification}
The most important problem the project addresses is that users
find it difficult to understand and assess complex legal documents in
the digital service industry. This is due to several
flaws in the system. Firstly, legal documents contain specialised legal language
and complex linguistic structures that are difficult to understand
for laymen. Secondly, the nature of agreements is such that it discourages
complete reading, where selective attention (or complete avoidance)
takes place. Thirdly, a layman may not have
the necessary legal knowledge to understand the implications of the contract
such that they are unable to point out clauses that might have a
potential impact on their rights and obligations.

In addition, there are no easy to use tools for pointing out
unfair or risky contract terms. The current state has
low levels of interactivity, which might be detrimental to the
users, as they may not be able to pose questions about certain
contract terms. All these issues make the users not aware of the context, 
which makes them vulnerable to privacy
infringements, financial risks, and service termination with
too little understanding.

\section{Problem Formulation}
Building on the difficulties we have described here, this project
explains the difficulty in terms of designing and putting into
practice an intelligent system for automatically evaluating Terms and
Conditions documents and giving users easily understandable,
risk-aware insights. The suggested system must fulfil a number of
usability and technical requirements in order to fully accomplish this goal.

Working with documents in different formats, extracting useful
semantic information from legal text, identifying and classifying
clauses\index{Clause Detection} based on potential risk, and offering explainable reasoning\index{Explainable AI}
to support risk assessments are some of the essential system
requirements. To facilitate deeper user understanding, the system
must also allow interactive querying and offer a trade-off between
user privacy requirements and performance.

This problem was chosen because of its applicability,
interdisciplinary characteristics, and possible implications for
society. The objective of this project is to enhance user
participation in digital decision-making processes and eliminate the
complexity hurdle that comes with legal documents using modern NLP
approaches, knowledge graphs, and interaction mechanisms.
